
\chapter{Basics}
\label{chap:basics}
The following chapter introduces the foundational theory and simulation techniques that will be used throughout this work. As our study focuses on turbulent boundary layers \benennung{\acrfull{tbl}}, we provide a brief overview of the theoretical background to these flows. Next, we discuss the necessity of this work within the context of relevant literature focusing on curved boundary layers. The final section outlines the applied optimisation method and provides an overview of the tools used. \info{Absatz sprachlich verbessert}


\section{Wall-bounded turbulence}\label{sec:turbulenceTheory}

% - boundary layer develops const in flow direction
% - with increasing boundary layer thickness
% - wall shear stress not know a priori
% - behaviour in inner layer ess. same as in channel flow
% - equation for pressure gradient
% %--> U_0 = const pressure gradient = 0
% - definition of boundary layer thickness, displ. thickness, momentum thickness
% - corresponding Reynolds numbers
% - shape factor

The description of fluids is based on three physical principles. \unsure{Kommen die drei Prinzipien jetzt sofort?, oder bezieht sich das auch auf den Text hinter Eq. 2.1?}
The first one  states that a fluid can be treated as a continuum \citep{schlichting_boundary-layer_2017} \change{Hört sich an, wie ne alte Erkenntnis, dann gehört hier als erstes eine alte Originalquelle hin}.
In addition, we take the viscous property of fluids into account, i.e.~a Newtonian fluid relates shear stresses to the rate of strain \citep{tennekes_first_1972}
Finally, it is assumed that the law of mass conservation  holds.
Using this, the \benennung{continuity equation} for viscous incompressible flows is given 
as\improve{Müssen die Variablen nicht eingeführt werden?, mind in einem Nomenklaturverzeochnis: $\rho: xxx$, gilt für die nächste Gleichung ebenfalls }

\begin{equation}
    \frac{\partial \rho}{\partial t} + \frac{\partial (\rho u_i)}{\partial x_i} = \frac{\partial u_i}{\partial x_i} = 0
\end{equation}

The second relevant equation describes the momentum conservation of a fluid.
These so-called \benennung{\acrfull{nse}} can be written as 

\begin{equation}
    \frac{\mathrm{D}u_i}{\mathrm{D}t} = \frac{\partial u_i}{\partial t} + u_j \frac{\partial u_i}{\partial x_j} = -\frac{1}{\rho} \frac{\partial p}{\partial x_i} + \nu \frac{\partial^2 u_i}{\partial x_j \partial x_j}
\end{equation}

These equations fully describe the motion of a fluid, since they are pplicable for most kind of flows (XXQUELLE).
Building on this,\unsure{However??}  the behaviour of fluids is often split into three groups with different characteristic properties:
These categories are defined as \emph{laminar},  \emph{turbulent} and  \emph{transitional}. 
If a flow acts in the laminar regime its streamlines do not cross and the flow is time-independent \citep{davidson_turbulence_2015}.
In the turbulent regime the behaviour of the flow is chaotic and time-dependent (XXQUELLLE).
It shows irregular time- and location-dependent random fluctuations \citep{pope_turbulent_2015}.
Consequently, in turbulent flows enhanced mixing in the fluid can be observed (XXQUELLE).
\improve{ein Halbsatz zu transitional}
To describe the behaviour of a flow in terms of laminar or turbulent behaviour the \benennung{Reynolds number} is introduced as 
\begin{equation}
    \mathrm{Re} = \frac{U L}{\nu}
\end{equation}
where $U$ and $L$ are reasonable velocity and length scales, and $\nu$ the viscosity of the fluid.
Flows with low Reynolds numbers show laminar flow behaviour.
With an increasing Reynolds number the fluid arrives first in the transitional region between laminar and turbulent until it reaches a fully turbulent state.
At high Reynolds numbers one can observe a separation of scales with different behaviours \citep{pope_turbulent_2015}.
One can split the turbulent motions into large- and small-scale motions.
The large scales are mostly controlled by the outer constraints of the flow, e.g., its geometry and the boundary conditions.
The behaviour of the small-scale motions, however, depends on the viscosity and the receiving energy from the large scales \citep{pope_turbulent_2015}, which makes their behaviour universal and independent of the specific flow case.
This universal behaviour can be measured with statistics over a long time averaging.

The most turbulent flows are bounded by at least one solid wall (XXQUELLE).
They are decomposed into \benennung{external} and \benennung{internal} flows. 
Internal flows are limited by outer walls like pipes or channels.
External flows refer to flows which occur around a solid body like an airfoil or a wall.
A characteristic feature of external flows is that the flow sufficiently from the body remains unaffected by its presence (XXQUELLE).
The flow behaviour in the immediate vicinity of the wall is in all these\unsure{Hier wurden mE keine Fälle aufgezählt: also besser "most cases"?} cases very similar due to the universal nature of small scales turbulences \citep{pope_turbulent_2015}XX umformulieren.
In the following section we will introduce a set of definitions to describe this near-wall behaviour.
is remarkably similar due to the universal nature of the near-wall small-scale turbulence \citep{pope_turbulent_2015}. 
%Therefore, it is useful to introduce a set of quantities that characterize this near-wall behaviour.

Since this work focuses on boundary layers developing along a wall, the definitions introduced in the following are presented specifically for this configuration.

\subsection{Turbulent boundary layer theory} \label{subsec:TBLTheory}

Boundary layers occur when fluids with non-zero viscosity stream over a wall that is not moving. 
In these setups the so called \benennung{no-slip} boundary condition impacts the flow behaviour significantly.
It states that the velocity components of the fluid in wall-normal and streamwise direction need to be zero directly at the wall since there is no relative movement between the wall and the fluid \citep{schlichting_boundary-layer_2017}. 
The velocity far away from the wall, the so called \benennung{free-stream velocity} $U_0$, is usually much higher than zero which implies the existence of a transition region where an acceleration of the fluid occurs (XXQUELLE).
This so-called \benennung{boundary layer} is placed close to the wall. 
Depending on the Reynolds number the thickness of the boundary layer varies in a way that the boundary layer decreases with increasing Reynolds number \cite{schlichting_boundary-layer_2017}. 
The flow in the boundary layer can behave laminar or turbulent, whereas in this thesis we will focus on \acrfull{tbl}.
The boundary layer grows in streamwise direction and can be described via multiple characteristic quantities. 
The introduction of these quantities requires a consistent coordinate system: 
We define $x$ as the streamwise direction with its corresponding velocity $u$, $y$ represents the wall-normal direction with velocity component $v$ and $z$ describes the spanwise direction with velocity component $w$.
In the free stream far away from the wall we assume $U_0$ being the dominant velocity component. 
For simplicity this free-stream velocity is often set to be 1. XX diesen satz vlt eher bei non-dim? XX
The streamwise velocity component inside the boundary layer is referred to as $u(x,y)$.
To describe the thickness of a boundary layer mainly three different parameters are used: 
the boundary-layer thickness $\delta_{99}(x)$, which describes the distance from the wall at which the velocity $u(x,y)$ derives $99\%$ of the free-stream velocity;
the displacement thickness $\delta^*$, which defines the displacement of an inviscid fluid to preserve the same mass flow;
and the momentum thickness $\theta$  which defines the displacement of an inviscid fluid to fulfil the conservation of momentum.
These quantities can be derived with the following equations
\begin{gather}\label{eq: tblThickness}
    u (x, \delta_{99}(x)) = 0.99 U_0 \\
    \delta^*(x) = \int_0^\infty \left( 1 - \frac{u(x,y)}{U_0} \right) \mathrm{d}y \hspace{0.8cm}
     \theta(x) = \int_0^\infty \frac{u(x,y)}{U_0} \left( 1 - \frac{u(x,y)}{U_0} \right) \mathrm{d}y
\end{gather}


In this thesis we use the momentum thickness $\theta$ to characterize the height of the boundary layer, since it can be measured more accurate than the boundary-layer thickness (XXQUELLE). 
To describe the boundary-layer with the Reynolds number, different Reynolds numbers based on the introduced displacement measures are defined as

\begin{gather}
    \mathrm{Re}_x = \frac{U_0 x}{\nu}, \quad
    \mathrm{Re}_{\delta_{99}} = \frac{U_0 \delta_{99}}{\nu}, \quad
    \mathrm{Re}_{\delta^*} = \frac{U_0 \delta^*}{\nu}, \quad
    \mathrm{Re}_\theta = \frac{U_0 \theta}{\nu}. \\
    \nu : \text{viscosity} \nonumber
\end{gather}

To measure velocities in a turbulent flow they are decomposed in a time averaged component $\langle \rangle$ and fluctuations $u_i'$. 

\begin{equation}
    u_i = \langle u_i \rangle + u_i' = U_i + u_i' \text{ with } \langle u_i'
\rangle = 0\end{equation}

represents the so-called \benennung{Reynolds decomposition}.
The application of the Reynolds decomposition on the \acrshort{nse} leads to the \acrfull{rans}. 
For the boundary layer the \acrshort{rans} equations reduces to 

\begin{gather}
    0 = \nu\,\frac{\mathrm{d}^2 U}{\mathrm{d}y^2} - \frac{\mathrm{d}}{\mathrm{d}y}\langle u'v'\rangle - \frac{1}{\rho}\frac{\partial \langle p\rangle}{\partial x} \\
    0 = -\frac{\mathrm{d}}{\mathrm{d}y}\langle v'v'\rangle - \frac{1}{\rho}\frac{\partial \langle p\rangle}{\partial y} 
\end{gather}

Through some rearrangements and boundary conditions one can derive the total shear stress $\tau$ and its derivative as

\begin{gather}
    \tau = \rho\nu\,\frac{\mathrm{d} U}{\mathrm{d}y} - \rho\langle u'v'\rangle \quad
    \frac{\mathrm{d}\tau}{\mathrm{d}y} = \frac{\mathrm{d}p_{\mathrm{w}}}{\mathrm{d}x}.
\end{gather}

To derive the shear stress at the wall one can apply the no-slip boundary condition which then results in
\begin{equation}
    \tau_w = \mu \left. \frac{\mathrm{d}U}{\mathrm{d}y} \right|_w 
\end{equation}

\comment{
%XX das hier weglassen statdessen mehr über non-dim und C_p/ C_f? 
From the \acrfull{nse} we can derive the boundary layer equations for a 2D boundary layer over a fixed flat wall for an incompressible fluid. 

\begin{align}
    \frac{\partial u}{\partial x} + \frac{\partial v}{\partial y} &= 0 \\
    u \frac{\partial u}{\partial x} + v \frac{\partial u}{\partial y} &= -\frac{1}{\rho} \frac{\partial p}{\partial x} + \nu \frac{\partial^2 u}{\partial y^2} \\
    0 &= - \frac{1}{\rho} \frac{\partial p}{\partial y} \label{eq:NSEy}
\end{align}

XX notwendig?


XX herleiten
With this equation we can define the wall shear stress $\tau_{w}$ as
\begin{equation}
    \tau_w := \mu \left. \frac{\mathrm{d}U}{\mathrm{d}y} \right|_w = -h \frac{\mathrm{d}p_w}{\mathrm{d}x}.
\end{equation}
}



\paragraph{Non-dimensionalisation}
To study the phenomenons in a boundary layer it is required to define so-called viscous scales since the ranges of dimensions and velocity variations are much smaller than in the free-stream \citep{pope_turbulent_2015}.
Based in the wall shear stress $\tau_w$, $\rho$ and $\delta_{99}$, the friction velocity $u_{\tau}$, the viscous length scale $\delta_{\nu}$ and the friction Reynolds number $\mathrm{Re}_{\tau}$ are introduced as
\begin{gather}
    u_{\tau} = \sqrt{\frac{\tau_w}{\rho}} \hspace{1.0cm}
    \delta_{\nu} = \frac{\nu}{u_{\tau}} \hspace{1.0cm}
    \mathrm{Re}_{\tau} = \frac{\delta_{99} u_\tau}{\nu}
\end{gather}

With these quantities the original length and velocity scales can be non-dimensionalised in wall units as \citep{moin_fundamentals_2024}
\begin{equation}
    y^+ = \frac{y}{\delta_v} = \frac{y u_\tau}{\nu}, \qquad
    u^+ = \frac{U}{u_\tau}.
\end{equation}

The advantage of these non-dimensional scales is that they are universal for all wall-bounded flows what enables the comparison between flows with different physical dimensions. XX unbedingt QUELLE XX
The non-dimensionalisation enables a detailed investigation of the velocity profile inside the \acrshort{tbl}.
Prandtl (XXQUELLE) registered that the \acrshort{tbl} contains an inner part where the mean velocity profiles is independent of the free-stream quantities $U_0$ and $\delta_{99}$ such that the mean velocity profile can be fully described by the plus units.
In the region closest to the wall – the \benennung{viscous sublayer} – $U^+$ is a linear function of $y^+$. 
With increasing wall-distance the velocity profile transitions such that it follows a logarithmic function.
This so-called \benennung{log law} is defined as 
\begin{gather}
    U^+ = \frac{1}{\kappa}\ln y^+ + B \\
    \text{with } \kappa = 0.41 \quad B = 5.2
\end{gather}

The values for the constants $\kappa$ and $B$ vary throughout the literature.
In the context of this work the values introduced by \citet{pope_turbulent_2015} are used as reference.
Typical velocity profiles from different turbulent flows are depicted in XX.
It can be seen that all profiles follow the described functions.

XX hier noch bild mit $U^+/y^+$


\subsubsection{Reynolds stresses}
\comment{
To measure velocities in a turbulent flow they are decomposed in a time averaged component $\langle \rangle$ and fluctuations $u_i'$. 

\begin{equation}
    u_i = \langle u_i \rangle + u_i' = U_i + u_i' \text{ with } \langle u_i'
\rangle = 0\end{equation}

represents the so-called \benennung{Reynolds decomposition}.
The application of the Reynolds decomposition on the \acrshort{nse} leads to the \acrfull{rans}. 

\comment{
\begin{equation}
    \frac{\partial \langle u_i \rangle}{\partial t} + \frac{\partial \langle u_i \rangle \langle u_j \rangle}{\partial x_j} = -\frac{1}{\rho} \frac{\partial \langle p \rangle}{\partial x_i} + \frac{1}{\rho} \frac{\partial}{\partial x_j} \left(  \mu \left( \frac{\partial \langle u_i \rangle}{\partial x_j} + \frac{\partial \langle u_j \rangle}{\partial x_i} \right) + \langle u_i' u_j'\rangle \right)
\end{equation}
}
}

As already introduced the application of the Reynolds decomposition on the \acrshort{nse} result in the \acrshort{rans}.
From this one can extract one term as the \benennung{Reynolds stresses}.
For turbulent flows the Reynolds stresses are important measures, since the transport of mean momentum is based on them \citep{moin_fundamentals_2024}. 
The tensor is defined as 
\begin{equation}
    R_{ij} =  \langle u_i' u_j'\rangle.
\end{equation}
The inner part $u_i' u_j'$ of the tensor elements represents the product of the velocity fluctuations.
To get the Reynolds stresses, the mean $\langle \rangle$ of this product is taken.
Although the mean of one fluctuation is zero, the product of two fluctuation components is not necessarily zero. 
The diagonal entries therefore, contain the normal stresses $\langle u' u'\rangle, \langle v' v'\rangle, \langle w' w'\rangle$.
The remaining entries contain the shear stresses.
Here the only non-zero entry is $\langle u' v'\rangle$ \citep{schlichting_boundary-layer_2017}.
The issue about this tensor is, that we cannot directly compute it.
With the continuity equation and the \acrshort{nse} in three directions we have four equations to derive the pressure and the three velocity components.
With the Reynolds stress tensor we have six additional terms we have to solve.
This problem is called the \benennung{closure problem} \citep{wilcox_turbulence_2006}.
To solve this, there are different approaches how to derive more equations such that all terms can be computed which is described in detail in fundamental books like \citet{wilcox_turbulence_2006,pope_turbulent_2015} XXnoch eines.
One can derive the transport equations for the Reynolds stresses by XXX \citep{mansour_reynolds-stress_1988}
For incompressible flows the non-dimensionalized transport equations can be stated as 
\begin{gather}
    0 = -\frac{\overline{\mathrm{D}}}{\overline{\mathrm{D}}t} \langle u_i u_j \rangle - T_{ij} + D_{ij} + \mathcal{P}_{ij} + \Pi_{ij} - \varepsilon_{ij}, \\[4pt]
    \begin{aligned}
        \mathcal{P}_{ij} &\colon \text{production rate,} & \quad \varepsilon_{ij} &\colon \text{dissipation rate,} \\
        T_{ij} &\colon \text{turbulent transport rate,} & \quad D_{ij} &\colon \text{viscous diffusion rate,} \\
        \Pi_{ij} &\colon \text{velocity pressure-gradient term,} & \quad \frac{\overline{\mathrm{D}}}{\overline{\mathrm{D}}t} \langle u_i u_j \rangle &\colon \text{convection}
    \end{aligned} \notag
\end{gather}
\comment{
\begin{equation}
\begin{aligned}
    0 = &-\left( \langle U \rangle \frac{\partial k}{\partial x} + \langle V \rangle \frac{\partial k}{\partial y} \right) + \mathcal{P} - \tilde{\varepsilon} + \nu \frac{\partial^2 k}{\partial y^2} \\
    &- \frac{\partial}{\partial y} \left\langle \frac{1}{2} v \boldsymbol{u} \cdot \boldsymbol{u} \right\rangle - \frac{1}{\rho} \frac{\partial}{\partial y} \langle v p' \rangle.
\end{aligned}
\end{equation}
\begin{equation}
    0 = -\frac{\overline{\mathrm{D}}}{\overline{\mathrm{D}}t} \langle u_i u_j \rangle - T_{ij} + D_{ij} + \mathcal{P}_{ij} + \Pi_{ij} - \varepsilon_{ij},
\end{equation}}


\subsection{Clauser parameter in turbulent boundary layers}

In boundary layers the free-stream (XX stimmt das?) pressure gradient plays a crucial role since it impacts the flow behaviour significantly (XXQUELLE).
From \autoref{eq:NSEy} we see that the pressure for a \acrshort{tbl} along a flat wall is constant in wall-normal direction.
In streamwise direction however it not necessarily has to be constant.
With Bernoulli's equation the free-stream pressure is connected to the free-stream velocity, such that the gradient in streamwise direction can be written like
\begin{equation}\label{eq:bernoulli}
    -\frac{\mathrm{d}p_0}{\mathrm{d}x} = \rho\, U_0 \frac{\mathrm{d}U_0}{\mathrm{d}x}.
\end{equation}
This equation shows that an acceleration of the flow effects a negative pressure gradient, which is also called \benennung{\acrfull{fpg}}. 
In contrast, an increase of the velocity leads to a positive or \benennung{\acrfull{apg}}. 
\acrshort{tbl} with a constant free-stream velocity therefore have no pressure gradient.
This \benennung{\acrfull{zpg}} configuration is well studied, since in plane \acrshort{tbl} it is the natural configuration (XXQUELLE).
However, it is possible to apply boundary conditions such that an adverse or favourable pressure gradient results.
To measure the pressure gradient a common used quantity is the so-called \benennung{Clauser parameter $\beta(s)$} which is defined as
\begin{equation}\label{eq:beta}
    \beta(s) = \frac{\delta^*(s)}{\tau_w(s)} \frac{\mathrm{d}p_w}{\mathrm{d}s},
\end{equation}

This non-dimensionalized parameter defines the ratio between the streamwise pressure force on the boundary layer and the wall shear stress \citep{clauser_turbulent_1954}.
In the literature this parameter is widely used to categorize the flow according to the introduced pressure regimes:

\begin{equation}
    \beta =
    \begin{cases}
        < 0, & \text{favourable pressure gradient (FPG)}, \\
        = 0, & \text{zero pressure gradient (ZPG)}, \\
        > 0, & \text{adverse pressure gradient (APG)}.
    \end{cases}
\end{equation}




\section{curved TBL}

Most research about \acrshort{tbl} focusses on \acrshort{tbl} across straight walls. (XXQUELLE) 
However, there are many real-world applications where the wall is not fully straight but shows curved sections as well (XXQUELLE). 
The effects in \acrshort{tbl} along curved walls are still poorly studied (XXQUELLE) since it covers some difficulties. 
Due to the curved shape a pressure gradient in streamwise direction builds up naturally which makes it hard to separate the source of observed physical effects \citep{so_experimental_nodate}. 
As \citet{so_experiment_1973} observed in their experiments without pressure adjustment the pressure decreases slow near the beginning of the curvature and then decreases fast once the curvature starts. 
Since these pressure variations have a high impact on the turbulent structures, an exclusive study of curvature effects requires a compensation of these pressure variations (XXQUELLE). 
Therefore, it is possible to study curved \acrshort{tbl} under zero-pressure gradient conditions (see XXX), or to impose a defined pressure gradient (see XXX). 
In experimental studies a defined pressure distribution is achieved by adjusting the position of the upper wall \citep{so_experimental_nodate, johnston_experimental_nodate, gillis_turbulent_1983,schwarz_convex_1996}. 
In simulations 
From these studies some general statements about the effects of curvature on the \acrshort{tbl} can be summarized.
First a boundary layer along a convex wall with \acrshort{zpg} shows smaller growth rate compared to a straight developing \acrfull{tbl} \citep{so_experimental_nodate}, \citep{meroney_turbulent_1975}. 
This phenomenon is caused by a pressure gradient in centripetal direction. 
It balances the centrifugal force of the fluid elements. Through the centrifugal force fluid particles with high velocities will be pushed to a new outer position. 
The pressure gradient forces the fluid particle in the opposite direction \citep{so_experimental_nodate}, \citep{meroney_turbulent_1975}.
\citet{gillis_turbulent_1983} observed a decrease in the turbulent scales as consequence of a convex curvature. 
In addition, they found that the law of the wall coincides for the curved wall with the same parameters as for flat wall measurements. 
Besides experimental investigations, some numerical studies about the characteristics of curved \acrshort{tbl} are presented.

Similar as in the experimental investigations a modification of the setup is necessary to compensate the streamwise pressure gradient \citet{appelbaum_systematic_2025}, XX. 
One possible approach to do so is the setup of a curved and a flat \acrshort{tbl} simulation and match their pressure distributions.
The naturally occurring pressure gradient in the curved domain is then imposed as boundary condition on the flat \acrshort{tbl} simulation.
With the flat plate as reference case one can thus isolate the effects of curvature.
This approach is used by \citet{spalart_direct_2024,pargal_adverse-pressure-gradient_2022}.
XXX hier bobke and schlatter einfügen
A similar approach \citet{bobke_history_2017} followed in their study about the effect of history on \acrshort{apg} \acrshort{tbl}.
To do this the $\beta(x)$ distribution from a wing \acrshort{dns} was copied on a flat \acrshort{tbl} through the imposition of corresponding velocity profiles along the top wall.
They studied various distributions and showed that even with matching $\beta(x)$ and $Re_{\tau}$ the geometry of the flow impacts the flow \citep{bobke_history_2017}.
Another possibility of course is the study of curved setups without manipulation of the pressure as done by \citet{vinuesa_pressure-gradient_2017,tanarro_effect_2020}, and XX.
\citet{vinuesa_pressure-gradient_2017} analysed data from flows along a wing which were derived by \citet{hosseini_direct_2016}.
They studied the occurring pressure field and found that even small variations of $\beta$ impacted the development of the boundary layer compared to \acrshort{zpg} \acrshort{tbl} \citet{vinuesa_pressure-gradient_2017}.
\citet{tanarro_effect_2020} extended these studies and analysed the different effects of local and integrated pressure gradient distributions.
However, these studies mainly focus on the effect of pressure on the flow characteristics refereeing to the argument that the pressure driven effects superpose the curvature effects \citep{patel_longitudinal_1997}.
\citet{yang_assessment_2016} studied a quite similar setup compared to the one used in the project to investigate curvature effects in \acrshort{tbl}. 
Their domain has a straight inlet, then a 60° circular bend section followed by a straight tail.
The \acrshort{rans} simulations they performed showed in comparison with \acrshort{dns} results, that accuracy of the tested turbulence models in insufficient in cases with large curvature \citep{yang_assessment_2016}.
In the study of \citep{appelbaum_systematic_2025} a concave \acrshort{tbl} was investigated. 
To achieve a \acrshort{zpg} distribution the pressure profile from the top wall of a plane \acrshort{tbl} was imposed on the top wall of the curved domain.
Since the pressure is the result of the flow behaviour it cannot be set directly an implicit approach was necessary to reach the desired pressure distribution \citep{appelbaum_systematic_2025}.
This overview of current research about curved \acrshort{tbl} shows that there is an interest in investigation of curvature effects on \acrshort{tbl}.
However, the separation of pressure and curvature effects is still an issue (XXQUELLE).
Many investigations focus on rebuilding the pressure distribution of a curved domain on a flat domain to get some understanding.
These studies already create some knowledge about curved \acrshort{tbl}.
However, the effects of curvature and pressure gradient still interfere with each other.
The study of the isolated curvature in simulations is rare since the realization of the \acrshort{zpg} is not straight-forward.
The purpose of this project is to implement and test a new workflow for setting up a curved \acrshort{zpg} \acrshort{tbl}.
The project is based on the previous work from \citet{morita_applying_2022} who introduced an algorithm to realize desired $\beta$-distributions.
The capability of the algorithm on curved domains was tested by \citet{ganagdhar_bayesian_2025}.
However, the results are only partly transferrable to the investigations from this study since the curvature studied by \citet{ganagdhar_bayesian_2025} had a much lower $\delta/R$ ratio. 
A detailed discussion about the transferability of other results is done in chapter XX.
The algorithm performs a shape optimization of the upper wall in an \name{OpenFOAM} simulation trying to fit a prescribed $\beta$ distribution at the lower wall.
Since we wanted to perform \acrshort{dns} simulations the result of the optimization need to be transferred.
With the optimized \acrshort{dns} domain we performed simulations with different curvature configurations.
The used approach is based on \citet{schlatter_leveraging_nodate} and will be presented in detail in \autoref{chap:methods}.

\comment{
To design a \acrshort{dns} domain with \acrshort{zpg} distribution the velocities along the upper wall will be manipulated.
The velocity components itself are extracted from optimized \acrshort{rans} simulations.
These \acrshort{rans} simulations are integrated in an optimization algorithm where the domain shape is manipulated to derive a desired $\beta$ distribution \citet{schlatter_leveraging_nodate}.
The optimization algorithm itself was introduced by \citet{morita_applying_2022}.
XX hier noch absatz 5.1 aus BO paper integrieren als motivation
}


\section{Simulation techniques} 
\label{sec:SimBasics}

For the simulation of turbulent flows different methods can be used. 
The decision about the used approach depends on the characteristics of the flow and the requirements for the solution.
One straight-forward approach is the \acrfull{dns}.
This method requires solving of the \acrshort{nse} to compute $U(x,t)$.
\acrshort{dns} is the computational most expensive approach since all length and time scales are resolved \citep{pope_turbulent_2015} which on the other hand gives the most accurate results.
Especially if the small turbulent scales are of interest, the usage of \acrshort{dns} is inevitable.
Because of the high computational costs \acrshort{dns} are restricted to low to moderate Reynolds numbers and mostly used for research.
The simulation method of \acrshort{dns} is mainly used within this project.
Additional simulations were performed with \acrfull{rans} simulations.
In this approach the \acrshort{rans} are solved instead of the full \acrshort{nse}.
As a consequence only the mean velocity field is calculated.
The effect of turbulence is summarized in the Reynolds stress tensor.
To avoid the computation of this tensor multiple approaches were developed to model it.
Since the models have different characteristics and accuracies the model choice can have significant impact on the simulation (XX QUELLE).
These simulations are computational much easier to realize.
If the exact behaviour of the turbulent scales is not important but only is interested in the effect of the turbulence on the flow \acrshort{rans} simulations build a reasonable compromise (XXQUELLE).
Therefore, this simulation method finds application in multiple industrial environments.


\section{optimization??}

\section{Tools}

XX das hier vlt in Basics?? XX

To work on this project multiple tools were used. Most of the fluid-flow simulations (korrekte bezeichnung??) were performed with DNS (see XX \autoref{sec:SimBasics}). To perform these simulations the spectral element code \name{neko} was used. Blabla über neko
When running simulations with \name{neko} it is possible to take statistics of the simulated flow during run time. Multiple settings about how the statistics should be recorded can be defined in the case file of the simulation. For example, it is possible to define one or two axis along which the statistics should be averaged. These statistics are used to extract and compute certain measures of the flow. To do this some post-processing of the recorded statistics is necessary. This is done in \name{python} where the package \name{pysemtools} is included. With this package data recorded in \name{neko} simulations can be extracted and interpolated from an SEM mesh into any arbitrary points (zu nah an github description). Then multiple statistical measures like the budget terms can be calculated. 
In addition to the DNS, some \acrshort{rans} simulations were performed as well. To run these simulations the open source software \name{OpenFOAM} is used. This software is able to perform \acrshort{rans} simulations of fluid flows with different turbulence models. To extract velocity fields etc from the simulation \name{Python} is used. 
All simulations are performed on the cluster of the RRZE of the FAU. Blabla über cluster
